\documentclass[11pt]{article}
\usepackage[utf8]{inputenc}
\usepackage{amsmath, amssymb}
\usepackage{geometry}
\geometry{a4paper, margin=1in}

\title{\textbf{Lecture 19: N-Variate Joint PDF/PMF/CDF and Jointly Gaussian Random Variables}}
\date{}

\begin{document}

\maketitle

\textit{(Primary Exam Scribe based strictly on Lecture PDF)}

\section*{1. Joint Distribution of N Random Variables}

\subsection*{1.1 Vector Representation}

For $N$ random variables $X_1, X_2, \dots, X_N$, the joint system is represented in \textbf{vector form}:

\[
\mathbf{X} =
\begin{bmatrix}
X_1 \\
X_2 \\
\vdots \\
X_N
\end{bmatrix}
\]

\begin{itemize}
\item This compact form allows us to treat multiple random variables as a \textbf{single entity}.
\item The \textbf{joint distribution} describes the \textbf{probability behavior of all variables together} (page 17).
\end{itemize}

\subsection*{1.2 Logical Transition: Bivariate $\rightarrow$ N-Variate}

\subsubsection*{Bivariate Case (2 variables)}

We describe:

\begin{itemize}
\item Joint PMF: $p_{X,Y}(x,y)$
\item Joint PDF: $f_{X,Y}(x,y)$
\item Joint CDF: $F_{X,Y}(x,y)$
\end{itemize}

\subsubsection*{Extension to N Variables}

The same idea scales \textbf{dimension-wise}:

\begin{itemize}
\item Instead of 2 variables $\rightarrow$ \textbf{N variables}
\item Instead of pair $(x,y)$ $\rightarrow$ \textbf{N-dimensional tuple} $(x_1, x_2, \dots, x_N)$
\end{itemize}

Thus, the structure generalizes as:

\subsubsection*{Joint PMF (Discrete Case)}

\[
p_{X_1, X_2, \dots, X_N}(x_1, x_2, \dots, x_N)
= \Pr(X_1 = x_1, X_2 = x_2, \dots, X_N = x_N)
\]

\subsubsection*{Joint PDF (Continuous Case)}

\[
f_{X_1, X_2, \dots, X_N}(x_1, x_2, \dots, x_N)
\]

\begin{itemize}
\item Represents density in \textbf{N-dimensional space}
\item Integrates over region $\rightarrow$ probability
\end{itemize}

\subsubsection*{Joint CDF}

\[
F_{X_1, \dots, X_N}(x_1, \dots, x_N)
= \Pr(X_1 \le x_1, \dots, X_N \le x_N)
\]

\subsubsection*{Key Insight (Scaling Logic)}

\begin{itemize}
\item The \textbf{definition does not change}, only:
\begin{itemize}
\item Number of variables increases
\item Domain becomes higher-dimensional
\end{itemize}
\item Conceptually: \textbf{Bivariate = 2D $\rightarrow$ N-variate = N-dimensional probability space}
\end{itemize}

\section*{2. Application Context: Large MIMO System}

From page 18:

\begin{itemize}
\item A wireless system has:
\begin{itemize}
\item $M$ transmit antennas
\item $N$ receive antennas
\end{itemize}
\end{itemize}

\subsection*{Interpretation}

\begin{itemize}
\item Each antenna produces a \textbf{random signal}
\item System forms an \textbf{N-variate random vector}
\item Joint distribution captures:
\begin{itemize}
\item Signal interactions
\item Correlations between antennas
\end{itemize}
\end{itemize}

\subsection*{Conclusion}

\[
\text{Number of random variables} \leftrightarrow \text{number of signals/antennas}
\]

\section*{3. Jointly Gaussian Random Variables}

\subsection*{3.1 Joint Gaussian PDF (Two Variables)}

For two variables $X$ and $Y$:

\[
f_{X,Y}(x,y)
=
\frac{1}{2\pi \sigma_X \sigma_Y \sqrt{1 - \rho^2}}
\exp \left(
-\frac{1}{2(1 - \rho^2)}
\left[
\left(\frac{x - \mu_X}{\sigma_X}\right)^2
+
\left(\frac{y - \mu_Y}{\sigma_Y}\right)^2
- 2\rho
\left(\frac{x - \mu_X}{\sigma_X}\right)
\left(\frac{y - \mu_Y}{\sigma_Y}\right)
\right]
\right)
\]

\subsection*{3.2 Parameter Interpretation (Strictly from Slides)}

\subsubsection*{Means}

\[
\mu_X, \mu_Y \rightarrow \text{centers of distributions}
\]

\subsubsection*{Standard Deviations}

\[
\sigma_X, \sigma_Y \rightarrow \text{spread of each variable}
\]

\subsubsection*{Correlation Coefficient}

\[
\rho = \frac{\text{cov}(X,Y)}{\sigma_X \sigma_Y}
\]

\begin{itemize}
\item Measures \textbf{linear dependence between X and Y} (page 7--8)
\end{itemize}

\subsection*{3.3 Structure of the Exponential Term}

The exponent contains three parts:

\subsubsection*{(1) Individual Contributions}

\[
\left(\frac{x - \mu_X}{\sigma_X}\right)^2, \quad
\left(\frac{y - \mu_Y}{\sigma_Y}\right)^2
\]

\begin{itemize}
\item Represent deviation of each variable independently
\end{itemize}

\subsubsection*{(2) Interaction Term}

\[
-2\rho
\left(\frac{x - \mu_X}{\sigma_X}\right)
\left(\frac{y - \mu_Y}{\sigma_Y}\right)
\]

\begin{itemize}
\item Captures \textbf{dependence between X and Y}
\end{itemize}

\subsubsection*{(3) Scaling Factor}

\[
\frac{1}{2(1 - \rho^2)}
\]

\begin{itemize}
\item Normalizes effect of correlation
\end{itemize}

\subsection*{3.4 Normalization Constant}

\[
\frac{1}{2\pi \sigma_X \sigma_Y \sqrt{1 - \rho^2}}
\]

\begin{itemize}
\item Ensures total probability integrates to \textbf{1}
\end{itemize}

\section*{4. Worked Example (From Slides)}

Given (page 4):

\[
\mu_X = 2, \quad \sigma_X^2 = 1
\]
\[
\mu_Y = -3, \quad \sigma_Y^2 = 4
\]
\[
\text{cov}(X,Y) = -1
\]

\subsection*{Step 1: Compute Standard Deviations}

\[
\sigma_X = 1, \quad \sigma_Y = 2
\]

\subsection*{Step 2: Correlation Coefficient}

\[
\rho_{XY} = \frac{-1}{(1)(2)} = -\frac{1}{2}
\]

\subsection*{Step 3: Substitute into Joint PDF}

Final result (page 10):

\[
f_{X,Y}(x,y)
=
\frac{1}{2\pi \sqrt{3}}
\exp \left(
-\frac{2}{3}
\left[
(x-2)^2
- \frac{(x-2)(y+3)}{2}
+ \frac{(y+3)^2}{4}
\right]
\right)
\]

\section*{5. Marginal Distributions}

From joint PDF (page 11--12):

\subsection*{Definition}

\[
f_X(x) = \int_{-\infty}^{\infty} f_{X,Y}(x,y)\,dy
\]

\[
f_Y(y) = \int_{-\infty}^{\infty} f_{X,Y}(x,y)\,dx
\]

\subsection*{Key Property (From Slides)}

If $(X, Y)$ are jointly Gaussian:

\[
X \sim \mathcal{N}(\mu_X, \sigma_X^2), \quad
Y \sim \mathcal{N}(\mu_Y, \sigma_Y^2)
\]

\section*{6. CDF and Q-Function}

\subsection*{Step 1: CDF Definition}

\[
F_X(x) = \Pr(X \le x)
\]

Thus:

\[
\Pr(X < 0) = F_X(0)
\]

\subsection*{Step 2: Standardization}

\[
Z = \frac{X - \mu_X}{\sigma_X}
\]

\[
\Pr(X < 0) = \Pr\left(Z < \frac{0 - 2}{1}\right)
= \Pr(Z < -2)
\]

\subsection*{Step 3: Symmetry}

\[
\Phi(-x) = 1 - \Phi(x) = \Pr(Z > x)
\]

\[
\Pr(Z < -2) = \Pr(Z > 2)
\]

\subsection*{Step 4: Q-Function}

\[
Q(x) = \Pr(Z > x)
\]

\subsection*{Final Answer}

\[
\Pr(X < 0) = Q(2)
\]

\end{document}
